\section{Chefs, Connoisseurs, and Consumers}

I've had restaurants recommended to me based on the quantity of food provided or the ``great'' salad bar, reasonable prices, or perhaps the convenient hours, and so on. These people are consumers of food; they expect the basics and are easily satisfied. Most restaurants depend on them for survival.

Connoisseurs, on the other hand, understand food preparation, the best combinations of foods and spices, and value the presentation almost as much as the flavor. Being more selective, they surpass the consumer. However, beyond the connoisseur is the chef, who prepares the meal. Such fare is not at all the most popular even if it is the best.

A similar ranking applies to weblogs. \textbf{Consumers} like to accumulate ideas because they are satisfying. The ideas may inspire them, confirm their own opinions, or add to the ideational superstructure that guides their lives. They are satisfied with the bon mot or the conclusion. How and why the conclusion is reached is of little import. Their favorite meal is the smorgasbord, a little of this and a little of that. Most weblogs are like that. They provide a series of ideas that are only loosely connected. They present a variety of concepts as long as they reduce to similar conclusions, even if the premises and the arguments are widely different.

The \textbf{connoisseur}, however, is more sophisticated. He has a reasonably developed worldview and both a wider and deeper background. He reads with understanding, he knows when his ideas are being confirmed and when they are challenged. For him, the argument means as much as the conclusion; in fact, he may even gain from following the discussion even when he rejects the conclusion.

Nevertheless, he is passive, since it is the \textbf{chef} who creates the argument. The reader is unaware of the outtakes, i.e., the text that didn't make the final cut. He usually misses the technique hidden within the text and may only be dimly aware of the true purposes of the chef. It is not to satisfy the connoisseur. Gornahoor is not interested in consumers or even in connoisseurs, but rather in training chefs.

\paragraph{The Medium is the Message}


\textbf{Marshall McCluhan} is famous for coining the expression that serves as the title for this section. In a work of art, for example, the unsophisticated consumer will take the ostensible object of the art as the important thing and ignore the framework. Hence they like images of kittens. A good critic, however, will focus on the technique; the kitten is of secondary interest. He will look at the technique, the use of color and shading, etc.

This past weekend, as we were listening to some music, I asked my older son, who is a musician, what exactly does he hear. At first, he didn't understand the question. I explained that most people just here pleasant sounds and they react to the beat. Then he understood. He told me he picks out the various voices, such as the bass line, the drum rhythms, and so on. He listens for the chord changes and can even tell what type of a chord is being played.

Recently, in a trad forum the topic of Dante came up. The moderator was nonplussed and objected that nothing could be learned from a mere work of fiction. All I can tell you is to abandon hope all ye who expect to restore tradition in that forum.

Clearly, Dante has a message, a quite sophisticated message in fact. Not only that, he actually explains the medium for those who might not figure it out for themselves. In the \emph{Convivio}, he explains the four levels of meaning that he inserted into his poetry: the literal, the allegorical, the moral, and the anagogical. To understand it as a mere work of fiction is to be stuck on the literal. It can also be read as theology, which many do. Ultimately, it is an initiatic work, i.e., it describes the stages on the spiritual path that are to be traversed in this life.

\paragraph{Being and Telling}


So a blog about Tradition must not be only about the message, that is, mere discourse in the manner of the professors. It itself must be written in a traditional way so that the medium itself becomes the message; it may even be the real message. Everything should have a purpose in relationship to the whole.

First of all, it must be organic. Hence, it is not systematic, yet it is not a mere smorgasbord of random ideas. Nevertheless, it is interconnected, and points in one post may be developed several months later, expressed in a different or deeper way.

It cannot be reduced to discursive thought, so there can be no explicit agenda or manifesto. However, there is still a point of view that is not subject to infinite debate. Therefore, it must be decisive when things deliquesce and solidity is restored.

It is dedicated to the Logos, even when that is not obvious and not incessantly proclaimed. To repeat, it is not talk about the Logos, but an attempt to manifest Him. Therefore, it strives to be logical and rational, in a transcendent sense, and it is hierarchically ordered.

It operates on the principle of affinity in the confidence that the right people and like spirits will congregate for their own reasons. Hence, there is no attempt at evangelization or to convince anyone through argument.

Organic thinking is creative thinking. That is why we present only original material, with the intent to provide a novel insight whenever possible. We don't simply republish another man's work. We do, however, publish our own translations of works that we consider useful. Some readers are interested only in those.

We also strive to write from our own understanding, whatever level that may be. Cynics will readily agree.

\paragraph{Fire and Ice}


Think carefully before committing to being a connoisseur or a chef. One will have to learn to be an actor to get along in the world, since everything is chosen freely. Usually, such a man will be careful about revealing his views; I am fortunate to be able to play the ``crazy uncle''. Detachment, when revealed, can be mistaken as inhuman coldness or even sociopathy. People can experience the Ice but not always the Fire.


\flrightit{Posted on 2013-03-28 by Cologero}

\begin{center}* * *\end{center}

\begin{footnotesize}
\begin{sffamily}
\texttt{graham on 2013-03-28 at 22:08 said:}

I'm sorry to hear about the snag. Dante's `metaphysics of dueling' appear to be true, but one must still take in the church's view of duels, judicial and otherwise, which, from what I've read, has been unfavourable, though other ordeals have met with cautious approval. Nevertheless I will see if we cannot look more into his technique.

\hfill

\texttt{Cologero on 2013-03-28 at 23:27 said:}

Vico provides an interesting justification for the duel. At the time, society provided for social disputes but not for disputes between individuals of the same rank. It was presumed that noble men would not lie. Now the purpose of a judge is to adjudicate disputes; but if both sides are telling the truth, on what basis can the judge make a judgment? If the disputants cannot resolve the issue between themselves, then the duel is the solution. It does have a certain logic. Clearly, it is not the Hollywood depiction of the duel that can only see hot tempers or evil wills.

\hfill

\texttt{Saladin on 2013-03-31 at 11:45 said:}

Cologero said `For both of them, it was always ``politics first''.'

I don't know Maurras (and frankly don't care) but why would you think Evola put politics first? Certainly he was not a Sannyasin and he definitely valued politics and angaged in it, but as far as I know he put Transcendence first and was well aware to which domain politics belonged.

\hfill

\texttt{Cologero on 2020-03-28 at 08:01 said:}

Saladin, Evola falsely believed the priestly caste was subservient to royalty. There was no reference to ``transcendence'' which you needlessly injected into the discussion. The proper domain of the political --- i.e., temporal power --- is to be subservient to the Spiritual Authority. It is not obvious that Evola accepted that.

\end{sffamily}
\end{footnotesize}

